\documentclass[11pt]{article}
\usepackage[T1]{fontenc}\usepackage{lmodern}
\usepackage[margin=1in]{geometry}
\usepackage{amsmath,amssymb,amsthm,graphicx,float,flafter,hyperref}
\hypersetup{colorlinks=true,linkcolor=blue,citecolor=blue,urlcolor=blue,pdftitle={Golden Coordinates at Level Five},pdfauthor={Akari Hayami (Jian-Yu Huang)}}
\setlength{\emergencystretch}{2em}
\newtheorem{theorem}{Theorem}[section]\newtheorem{proposition}[theorem]{Proposition}\newtheorem{corollary}[theorem]{Corollary}
\newcommand{\Q}{\mathbb Q}\newcommand{\R}{\mathbb R}\newcommand{\ii}{\mathrm i}\newcommand{\e}{\mathrm e}\newcommand{\Fcal}{\mathcal F}\newcommand{\Tcal}{\mathcal T}
\newcommand{\BodyFigure}[3]{\begin{figure}[tbp]\centering
\includegraphics[width=\linewidth]{companion/body/pdf/body_#1.pdf}
\caption{#2}\label{#3}\end{figure}}
\newcommand{\BodyFigureHere}[3]{\begin{figure}[H]\centering
\includegraphics[width=\linewidth]{companion/body/pdf/body_#1.pdf}
\caption{#2}\label{#3}\end{figure}}
\renewcommand{\topfraction}{.9}\renewcommand{\bottomfraction}{.75}
\renewcommand{\textfraction}{.08}\renewcommand{\floatpagefraction}{.8}
\raggedbottom
\title{Golden Coordinates at Level Five: Ruled-Surface Rank Loss, CM Values, and Fourier Self-Duality}
\author{Akari Hayami (Jian-Yu Huang)\\{\small\href{mailto:akari.h.1.9984@gmail.com}{\nolinkurl{akari.h.1.9984@gmail.com}}}}\date{20 October 2026\\Portrait revision 0.2}
\begin{document}\maketitle
\begin{abstract}
The interior rank loss of an orthogonal-circle ruled surface selects
$c_0=(3-\sqrt5)/2=\phi^{-2}$. We give explicit maps connecting this coordinate
to level-five arithmetic and Fourier symmetry. At
$\tau_{163}=(1+\ii\sqrt{163})/2$, an exact factorization of Klein's equation
recovers $\sqrt5$ and hence $c_0$ from $v=r(\tau_{163})^5$. A square witness
identifies $\Q(v)=\Q(v-v^{-1})(\sqrt{163(1+c_0)})$; the radical is a
pentagon side length at circumradius $\sqrt{163}$. At the distinct point $i$,
Rogers--Ramanujan inversion is conjugate by coordinate exchange to the odd
period-five sine-Fourier involution. Its two eigenlines characterize the two
signed completed functional equations by Poisson--Mellin analysis. Exact
arithmetic certificates are kept separate from the analytic arguments.
\end{abstract}

\section{Introduction and principal assertions}\label{sec:intro}
Hermite's 1859 discussion leads, in modern notation, to the remarkably small
gap between $\e^{\pi\sqrt{163}}$ and an integer \cite[pp.~47--48]{Hermite}.  For
$\tau_{163}=(1+\ii\sqrt{163})/2$ and
$q=\e^{2\pi\ii\tau_{163}}=-\e^{-\pi\sqrt{163}}$, the classical
$j$-expansion $q^{-1}+744+196884q+\cdots$ and
$j(\tau_{163})=-640320^3$ explain its small gap \cite{Duke}.
Meanwhile $P(c)=c^2-3c+1$ has discriminant $5$, so its roots lie in
$\Q(\sqrt5)=\Q(\zeta_5+\zeta_5^{-1})$.
Indeed $\zeta_5+\zeta_5^{-1}=2\cos(2\pi/5)=(\sqrt5-1)/2$ has minimal
polynomial $X^2+X-1$, identifying the real quadratic cyclotomic subfield.  The aim is to identify maps that recover this scalar from a CM value and
carry it into Fourier symmetry.

Throughout, put $\phi=(1+\sqrt5)/2$, $\gamma=\phi^{-1}$,
$c_0=\phi^{-2}=\gamma^2$, and $u_0=\gamma/2$.
The arithmetic connection is explicit: Theorem~\ref{recovery} recovers
$\sqrt5$ as a rational function of the actual CM value $v=r(\tau_{163})^5$,
where $r$ is the Rogers--Ramanujan function defined in \eqref{rr}.
Corollary~\ref{equality} then places the length
$\delta=\sqrt{163(1+c_0)}=2\sqrt{163}\sin(\pi/5)$ inside $\Q(v)$.
The companion length $\phi\delta$ is the diagonal of the same regular
pentagon. Thus its squared side-to-diagonal ratio is precisely the surface
coordinate $c_0$.

The second connection identifies two actions on coordinates. The matrices in
Section~\ref{sec:involution} satisfy $A=JMJ$: normalized modular inversion
becomes the odd five-periodic sine-Fourier involution after exchanging the two
coordinates. Proposition~\ref{kappa} identifies its positive eigenline through
$\kappa=r(i)$. Theorems~\ref{poisson} and~\ref{classification} carry this
finite symmetry to an entire completed function and characterize both signed
functional equations. The CM value and modular inversion are classical inputs
from Duke; the rational recovery, square witness and coordinate maps are
explicit here, and the analytic step uses Poisson summation. Figure~\ref{map}
separates the arithmetic role of $\tau_{163}$ from the fixed-point role of $i$.

\BodyFigure{01}{The two CM arguments enter at different stages. The value at
$\tau_{163}$ recovers the golden scalar; the fixed point $i$ selects the
positive Fourier eigenline. The arrows summarize the explicit maps proved
in the text.}{map}

\section{The surface and its golden coordinate}\label{sec:surface}
Following the paired-circle construction \cite[Definition~1.1]{HMR}, use
$-\pi/2<t<\pi/2$ and $0\le u\le1$. The surface angle is $t$ and the
ruling coordinate is $u$:
\begin{equation}\label{surface}
 S(t,u)=(c+2u(1-c),(1-u)\sin t,u\sqrt{c(2-c)}),\qquad c=\cos t.
\end{equation}
The paired endpoints are $C_1(t)=(c,\sin t,0)$ and
$C_2(t)=(2-c,0,\sqrt{c(2-c)})$: points on unit-circle arcs in perpendicular planes,
with centers $(0,0,0)$ and $(1,0,0)$.  Thus
$S(t,u)=(1-u)C_1(t)+uC_2(t)$ for $0\le u\le1$.
For fixed $t$, varying $u$ traces a ruling segment.  Rank loss means that
its two parameter-direction tangent vectors become linearly dependent.  In $0<c<1$, the third component of
$S_t\times S_u$ is $(1-c)(1-c-2u)$, hence rank loss forces $u=(1-c)/2$.
The remaining components are
\[
-\frac{(1+c)P(c)}{2\sqrt{c(2-c)}},\qquad
-\frac{\sin t\,P(c)}{\sqrt{c(2-c)}}.
\]
The prefactor of the first remaining component is nonzero in $0<c<1$.
Consequently, in this interior range,
\[
\operatorname{rank}(dS)<2
\quad\Longleftrightarrow\quad u=(1-c)/2\ \text{and}\ P(c)=0.
\]
The unique root of $P$ in $(0,1)$ therefore gives
\begin{equation}\label{golden}
c_0=\frac{3-\sqrt5}{2},\qquad\gamma=1-c_0=\phi^{-1},\qquad
u_0=\gamma/2=\cos(2\pi/5).
\end{equation}
The corresponding parameter points are
$(t,u)=(\pm\arccos c_0,u_0)$; we call them the reflected lateral pair.
Figure~\ref{surfacefig} displays the surface and the two conditions for rank
loss. This local calculation is also the interior part of the published
five-singularity classification \cite[Theorem~2.2]{HMR}.

\BodyFigure{02}{The actual ruled patch and its rank-loss diagnostic. On the
positive-sine parameter branch, the dashed line $u=(1-c)/2$ is necessary;
rank loss occurs at its marked point $(c_0,u_0)$, where $P(c_0)=0$.
The color scale is $\log_{10}(\|S_t\times S_u\|+10^{-8})$.
The surface panel marks both reflected parameter points.}{surfacefig}

The projection $u_0=\cos72^\circ$ belongs to a unit pentagon, whose squared
side-to-diagonal ratio is $c_0$. A second Euclidean realization uses the
unit-radius icosahedron with vertices
\[
\frac{(0,\pm1,\pm\phi),\ (\pm1,\pm\phi,0),\ (\pm\phi,0,\pm1)}
{\sqrt{1+\phi^2}}.
\]
Its distinct dot products are $1/\sqrt5,-1/\sqrt5,-1$, with pair counts
$30,30,6$. The squared short-to-long chord ratio is
$(2-2/\sqrt5)/(2+2/\sqrt5)=c_0$. Supplementary plates S04 and S12 show these
models; the full surface and parameter diagnostics are in S02--S03.

\section{CM field recovery}\label{sec:cm}
For $\Im\tau>0$, use the Rogers--Ramanujan product
\cite[Eq.~(4.1)]{Duke}
\begin{equation}\label{rr}
r(\tau)=q^{1/5}\prod_{n\ge1}\frac{(1-q^{5n-4})(1-q^{5n-1})}{(1-q^{5n-3})(1-q^{5n-2})},
\quad q=\e^{2\pi\ii\tau},\quad q^{1/5}=\e^{2\pi\ii\tau/5}.
\end{equation}
The two arguments have the following roles:
\begin{center}
\begin{tabular}{@{}llll@{}}
\hline
Argument & Nome $q$ & $j$-value & Role\\ \hline
$i$ & $\e^{-2\pi}$ & $1728$ & inversion fixed point\\
$\tau_{163}$ & $-\e^{-\pi\sqrt{163}}$ & $-640320^3$ & arithmetic recovery\\ \hline
\end{tabular}
\end{center}
Set $v=r(\tau_{163})^5$. The classical CM value above and Klein's
level-five relation (as recorded by Duke \cite[Eq.~(2.5)]{Duke}) give
\begin{equation}\label{klein}
P_j(v)=(v^4-228v^3+494v^2+228v+1)^3+jv(v^2+11v-1)^5=0.
\end{equation}
Let $A_{163},B_{163}\in\mathbb Z[X]$ be the explicit polynomials
in Appendix~\ref{certificate}.
\begin{theorem}\label{recovery}
For $v=r(\tau_{163})^5$, $\sqrt5\in\Q(v)$ and
\[
\begin{aligned}
\sqrt5&=-\frac{A_{163}(v)}{B_{163}(v)},&
\phi&=\frac{B_{163}(v)-A_{163}(v)}{2B_{163}(v)},\\
\gamma&=-\frac{A_{163}(v)+B_{163}(v)}{2B_{163}(v)},&
c_0&=\frac{3B_{163}(v)+A_{163}(v)}{2B_{163}(v)}.
\end{aligned}
\]
\end{theorem}
\begin{proof}
Appendix~\ref{certificate} gives the exact identity
$P_{-640320^3}(X)=A_{163}(X)^2-5B_{163}(X)^2$ with coprime integer
polynomials.  If $B_{163}(v)=0$, then $A_{163}(v)=0$, contrary to Bezout.
Thus $(A_{163}(v)/B_{163}(v))^2=5$.  The rational interval certificate gives
$A_{163}(v)<0<B_{163}(v)$, selecting $-\sqrt5$; the formulas follow.
\end{proof}
The square witness in Appendix~\ref{certificate} strengthens this recovery:
with $U=v-v^{-1}$, it proves
\[
\Q(v)=\Q(U)(\delta),\qquad \delta=2\sqrt{163}\sin(\pi/5).
\]
The pentagonal side length is therefore represented inside the field of the
selected CM value by an explicit formula, not just by its defining radical.

For the near-integer motivation, put $R=\e^{\pi\sqrt{163}}$ and
$N=640320^3+744=262537412640768744$. Substituting $q=-1/R$ in the classical
$j$-series gives
\[
N-R=\frac{196884}{R}-\frac{21493760}{R^2}
 +\frac{864299970}{R^3}+\cdots
 \approx 7.499274028\times10^{-13}.
\]
For these coefficients, $j=E_4^3/\Delta$, where
$E_4(q)=1+240\sum_{n\ge1}\sigma_3(n)q^n$,
$\Delta(q)=q\prod_{n\ge1}(1-q^n)^{24}$, and
$\sigma_3(n)=\sum_{d\mid n}d^3$.
Figure~\ref{near} shows successive corrections to this one fixed gap.
The complete CM comparison and numerical ladder are in S10--S11.

\BodyFigureHere{03}{Successive $j$-series corrections to the fixed gap $N-R$.
For correction count $m$, the residual is
$|N-R-\sum_{n=1}^{m}(-1)^{n+1}c_n/R^n|$, where
$c_1=196884$, $c_2=21493760$, $c_3=864299970$.
Values use 180 decimal digits and six corrections; the series coefficients
are obtained by exact expansion of $E_4^3/\Delta$.}{near}

\section{The common involution}\label{sec:involution}
Duke's inversion law \cite[Proposition~2, Eq.~(3.2)]{Duke} is
\[
r(-1/\tau)=\frac{\gamma-r(\tau)}{1+\gamma r(\tau)}.
\]
In homogeneous coordinates put
\[
M=\frac{1}{\sqrt{1+c_0}}\begin{pmatrix}-1&\gamma\\\gamma&1\end{pmatrix},
\qquad J=\begin{pmatrix}0&1\\1&0\end{pmatrix}.
\]
Taking the ratio of $M(z,1)^T$ gives the displayed law, and direct
multiplication, using $\gamma^2=c_0$, proves $M^2=I$.

Let $V_5$ be the real space of odd five-periodic sequences.  In residues
$1,2,3,4,5$ (with $5=0$), it consists of $(x,y,-y,-x,0)$.  Define
\[
(\Fcal a)_m=\frac{1}{\sqrt5}\sum_{n=1}^5a_n\e^{2\pi\ii mn/5},\qquad
\Tcal=-\ii\Fcal.
\]
Pairing $n$ and $5-n$ gives
\begin{equation}\label{A}
A=\frac{2}{\sqrt5}\begin{pmatrix}\sin(2\pi/5)&\sin(\pi/5)\\
\sin(\pi/5)&-\sin(2\pi/5)\end{pmatrix}=JMJ.
\end{equation}
Indeed $2\sin(2\pi/5)/\sqrt5=1/\sqrt{1+c_0}$ and
$2\sin(\pi/5)/\sqrt5=\gamma/\sqrt{1+c_0}$, by squared pentagon
identities.  Thus $A^2=I$.

\begin{proposition}\label{kappa}
With $\kappa=r(i)$,
\[
\kappa=\frac{\gamma}{1+\sqrt{1+c_0}},\qquad A(1,\kappa)^T=(1,\kappa)^T.
\]
For $\theta=\frac12\arctan\gamma$,
\[
\kappa=\tan\theta,\quad\cos2\theta=\frac{1}{\sqrt{1+c_0}},\quad
\sin2\theta=\frac{\gamma}{\sqrt{1+c_0}},
\]
and, separately, $\kappa=2\cos18^\circ-2\cos36^\circ$.  Finally
$u_0=\kappa/(1-\kappa^2)$ and
$c_0=(2\kappa/(1-\kappa^2))^2$.
\end{proposition}
\begin{proof}
The inversion law fixes $r(i)$; positivity of every product factor in
\eqref{rr} at $i$ selects the positive root of
$\gamma X^2+2X-\gamma=0$.  It gives the stated $\kappa$ and
$M(\kappa,1)^T=(\kappa,1)^T$, so $J$ transfers it explicitly to the
$+1$ eigenvector of $A$.  Since $\tan2\theta=\gamma$ and
$0<2\theta<\pi/2$, right-triangle and half-angle formulas prove the
three half-angle identities.  This axis angle is not $18^\circ$.
Separately $2\cos36^\circ=\phi$, $2\cos18^\circ=\sqrt{2+\phi}$;
squaring positive sides shows their difference has the same positive
quadratic root.  The last formulas are the tangent double-angle identity.
\end{proof}
The fixed-point polynomial also has a norm interpretation:
\[
X^4+2X^3-6X^2-2X+1=(X^2+2\phi X-1)(X^2-2\phi^{-1}X-1),
\]
the polynomial norm
$N_{\Q(\sqrt5)(X)/\Q(X)}(X^2+2\phi X-1)$.
Conjugation sends $\phi$ to $-\phi^{-1}$; multiplication using
$\phi-\phi^{-1}=1$ and $\phi\phi^{-1}=1$ proves the identity.
The fixed-point equation gives $\kappa^2+2\phi\kappa-1=0$,
so the first factor vanishes at $X=\kappa$.

Figure~\ref{exchange} shows the coordinate exchange and the positive
Fourier coefficient pattern. Supplementary plates S05--S07 give the full
reflection, direct-product and finite-DFT diagnostics.
\BodyFigureHere{04}{Coordinate exchange between the reflections $M$ and $A=JMJ$.
The paired off-axis inputs are $\mathbf z$ and $J\mathbf z$; their images
satisfy $A(J\mathbf z)=J(M\mathbf z)$. The solid axes are the positive
eigenlines, with directions $(\kappa,1)$ and $(1,\kappa)$ respectively;
the dashed axes have eigenvalue $-1$. The residue strip gives the normalized
positive Fourier eigenvector in order $1,2,3,4,5=0$.}{exchange}

\section{Poisson--Mellin bridge and classification}\label{sec:theta}
The Gaussian bridge turns finite Fourier symmetry into analytic symmetry.
To retain $t$ as a surface angle, use $x>0$ here:
\[
F_a(s)=\sum_{n\ge1}a_n n^{-s}\quad(\Re s>1),\qquad
\Theta_a(x)=\sum_{n\in\mathbb Z}na_n\e^{-\pi n^2x/5}.
\]
\begin{theorem}\label{poisson}
For $a\in V_5$, $a'=\Tcal a$,
\[
\Lambda_a(s)=\frac{1}{2}\int_0^\infty\Theta_a(x)x^{(s+1)/2}\frac{dx}{x}
=\left(\frac{5}{\pi}\right)^{(s+1)/2}\Gamma((s+1)/2)F_a(s)
\]
defines an entire function by its integral; the displayed Dirichlet-series
expression holds for $\Re s>1$.  Moreover $\Lambda_a(s)=\Lambda_{a'}(1-s)$.
\end{theorem}
\begin{proof}
With $\widehat f(\xi)=\int_\R f(y)\e^{-2\pi\ii y\xi}dy$,
put $f_x(y)=y\e^{-\pi xy^2/5}$.  Its transform is
\[\widehat f_x(\xi)=-\ii(5/x)^{3/2}\xi\e^{-5\pi\xi^2/x}.\]
Residue-class Poisson summation \cite[\S1.8(iv)]{DLMF} is
\[
\sum_{\ell\in\mathbb Z}f_x(j+5\ell)=\frac{1}{5}\sum_{k\in\mathbb Z}
\e^{2\pi\ii kj/5}\widehat f_x(k/5).
\]
After multiplying by $a_j$ and summing,
\begin{align*}
\Theta_a(x)&=\frac{1}{\sqrt5}\sum_k(\Fcal a)_k\widehat f_x(k/5)\\
&=\frac{1}{\sqrt5}\sum_k\ii a'_k[-\ii(5/x)^{3/2}(k/5)
\e^{-\pi k^2/(5x)}]=x^{-3/2}\Theta_{a'}(1/x).
\end{align*}
This shows the $1/5$, $1/\sqrt5$, $\ii(-\ii)$ cancellation, and
$x^{-3/2}$.  Bounded periodicity gives exponential decay at infinity;
the transformation gives $Cx^{-3/2}\e^{-c/x}$ at zero.  These envelopes,
also after multiplication by $|\log x|^r$, prove local uniform convergence
of the integral and all $s$ derivatives.  For $\Re s>1$, odd pairing and
termwise gamma integration give the displayed completion.  Substitution of
the theta transformation changes $x^{(s+1)/2}$ to $x^{(s-2)/2}$, and
$y=1/x$ gives $y^{(2-s)/2}dy/y$, proving the functional equation.
\end{proof}

\begin{theorem}\label{classification}
For nonzero $a\in V_5$ and $\epsilon\in\{+1,-1\}$,
$\Lambda_a(s)=\epsilon\Lambda_a(1-s)$ for all $s$ iff $\Tcal a=\epsilon a$.
The positive line is $\R(1,\kappa,-\kappa,-1,0)$, fixed by $a_1=1$;
the negative line is $\R(1,-1/\kappa,1/\kappa,-1,0)$, and its completed
value at $1/2$ is zero.
\end{theorem}
\begin{proof}
Root-of-unity orthogonality gives $\Fcal^2a=a_{-n}=-a$, hence $\Tcal^2=I$.
Solving the two eigenline equations of \eqref{A} gives the stated lines.
One implication follows from Theorem~\ref{poisson}.  Conversely let
$b=\Tcal a-\epsilon a$.  The assumed equation and that theorem make
$F_b$ vanish in a right half-plane.  If $m$ is its first nonzero coefficient,
then for real $\sigma\to\infty$,
\[
0=m^\sigma F_b(\sigma)=b_m+\sum_{n>m}b_n(m/n)^\sigma.
\]
The bounded tail tends to zero because
\[
\left|\sum_{n>m}b_n(m/n)^\sigma\right|
\le C m^\sigma\int_m^\infty x^{-\sigma}dx
=\frac{Cm}{\sigma-1}.
\]
Thus $b_m=0$, a contradiction.
Central vanishing is immediate for $\epsilon=-1$.
\end{proof}
For a general vector, put $a_\pm=(a\pm\Tcal a)/2$.  Its completion
is the sum of a symmetric and an antisymmetric part.  If both components
are nonzero, neither scalar-sign equation holds.  For example, the vector with $(x,y)=(1,0)$
is an eigenvector of neither sign.  The normalized $+1$ line is selected by
$a_1=1$ and gives the Davenport--Heilbronn sequence; its parameter is the
classical constant $\kappa=(\sqrt{10-2\sqrt5}-2)/(\sqrt5-1)$.

Here $a_6=1$ but $a_2a_3=-\kappa^2$, so the normalized $+1$ sequence is
nonmultiplicative.  The off-critical-line zero result is classical and
externally attributed \cite{Balanzario}; this is only a classification inside
$V_5$.

For the normalized positive vector
$a=(1,\kappa,-\kappa,-1,0)$, the paired theta transformation becomes
$\Theta_a(x)=x^{-3/2}\Theta_a(1/x)$. Supplementary plates S08--S09 show
independent numerical evaluations of this identity and the completed function.
On a fixed-height complex slice, evenness of the real part and oddness of the
imaginary part about $\Re s=1/2$ use both the functional equation and
$\Lambda_a(\overline s)=\overline{\Lambda_a(s)}$, since the coefficients
are real. The numerical methods and saved data are stated in the companion.

\section{Conclusion}
The surface calculation supplies the golden coordinate, the CM certificate
recovers it from the actual $\tau_{163}$ value, and $A=JMJ$ supplies the
analytic payoff at $i$.  The recovered golden scalar determines a reflection,
whose positive eigendirection determines the analytic symmetry.  The explicit
maps connect the geometric ratio, arithmetic field and functional equation
through transformations that can each be checked.  A precise remaining question is whether an
$N$-dependent ruled-surface family can be constructed to carry the odd-$N$
Fourier mechanism below.

\paragraph{Associated files.}
The accompanying source bundle contains the article, the visual companion
(plates S01--S12), the plotting data and scripts, and exact verification
programs. Its README gives replay commands and file hashes. The geometric
source used in \cite{HMR} is included with its published-source hash.

\appendix
\section{Exact certificate and square witness}\label{certificate}
Put $U=v-v^{-1}$ and
\begin{align*}
\alpha(U)&=U^3-p_1U^2+(3-p_2)U-(2p_1+p_3),\\
\beta(U)&=-d_1U^2-d_2U-(2d_1+d_3).
\end{align*}
The integer coefficients are
\[
\begin{array}{ll}
p_1=131268706320384342,&d_1=58705150130167150,\\
p_2=2896632145572621515,&d_2=1295413276662289800,\\
p_3=15717124384666951860,&d_3=7028911706986972100.
\end{array}
\]
The reconstruction polynomials are
\begin{align*}
A_{163}(X)&=X^6-p_1X^5-p_2X^4-p_3X^3+p_2X^2-p_1X-1,\\
B_{163}(X)&=-d_1X^5-d_2X^4-d_3X^3+d_2X^2-d_1X.
\end{align*}
Exact expansion gives
$P_{-640320^3}(X)=A_{163}(X)^2-5B_{163}(X)^2$ and
$\gcd(A_{163},B_{163})=1$, equivalently
$P_{-640320^3}(v)=v^6(\alpha(U)^2-5\beta(U)^2)$.

Define $G(q)$ as the positive product in \eqref{rr}, so $v=qG(q)^5$.
For $q=-\e^{-\pi\sqrt{163}}$ every factor is positive and the phase is
$r(\tau_{163})=\e^{\pi\ii/5}(-v)^{1/5}$.  The rational estimates
\[
3.14<\pi<22/7,\quad12.76<\sqrt{163}<12.77,\quad
\e^{40}>23\cdot10^{16},\quad\e^{41}<7\cdot10^{17},
\]
\[
|\log G(q)|\le \frac{|q|}{(1-|q|)^2}
\]
give $40<\pi\sqrt{163}<41$, $5|\log G(q)|<1/1000$, and
\[
-5\cdot10^{-18}<v<-10^{-18}.
\]
Exact interval Horner arithmetic gives $A_{163}(v)<0<B_{163}(v)$
throughout this rational interval; a Sturm count gives exactly one root of
$P_{-640320^3}(X)$ there. The source bundle replays these exact finite
computations and records separate 100/200-digit product diagnostics.

With $S=\sqrt5$, define
\begin{align*}
w(U)=\frac{1}{1574580000}[&(312119004989-139583862445S)U^2\\
&+(6858643496508-3067324783740S)U\\
&+37679427598544-16850735555920S].
\end{align*}
Exact reduction by $\alpha(U)+S\beta(U)$ and $S^2-5$ proves
\[
U^2+4=163(1+c_0)w(U)^2.
\]
On the selected branch $U>1$, the leading and linear coefficients are
negative and $w(1)<0$, hence $w'(U)<0$ and $w(U)<0$.
\begin{corollary}\label{equality}
For $\delta=\sqrt{163(1+c_0)}>0$, $\Q(v)=\Q(U)(\delta)$.
\end{corollary}
\begin{proof}
The coefficient formulas also give
$A_{163}(v)=v^3\alpha(U)$ and $B_{163}(v)=v^3\beta(U)$.
Hence $S=-\alpha(U)/\beta(U)\in\Q(U)$ and $w(U)\in\Q(U)$.
Since $U\in\Q(v)$, both also belong to $\Q(v)$.  Since
$2v-U=v+v^{-1}<0$ and $w(U)<0$,
$\delta=(2v-U)/w(U)\in\Q(v)$.  Conversely
$v=(U+\delta w(U))/2\in\Q(U)(\delta)$.
\end{proof}

The square-witness verifier in the source bundle replays the zero remainder
and the rational sign comparisons.
The selected radical has a direct pentagonal interpretation:
\[
\delta=2\sqrt{163}\sin(\pi/5),\qquad
\phi\delta=2\sqrt{163}\sin(2\pi/5).
\]
Indeed $1+c_0=(5-\sqrt5)/2=4\sin^2(\pi/5)$ and
$2\cos(\pi/5)=\phi$.  The two lengths are the side and diagonal
of a regular pentagon of circumradius $\sqrt{163}$; their squared ratio
is $c_0$.  Eliminating $S$ from $\delta^2=163(5-S)/2$ gives
$\delta^4-815\delta^2+132845=0$, a different polynomial from the
norm quartic for $\kappa$.  The cyclic quartic field of conductor $815$ is a
known family \cite{SpearmanWilliams}; the polynomial here records the selected
generator used above.

\section{Odd-\texorpdfstring{$N$}{N} extension}\label{oddN}
For odd $N>1$, let $V_N$ be real odd $N$-periodic sequences and define
\[
(\Fcal_Na)_m=N^{-1/2}\sum_{n=1}^Na_n\e^{2\pi\ii mn/N},\qquad
\Tcal_N=-\ii\Fcal_N.
\]
Its real sine matrix is
\[
\frac2{\sqrt N}\bigl(\sin(2\pi mn/N)\bigr)_{1\le m,n\le(N-1)/2}.
\]
Orthogonality gives $\Fcal_N^2a=a_{-n}=-a$, so $\Tcal_N^2=I$; no unique
positive line is claimed.  With
\[
\Theta_{a,N}(x)=\sum_{n\in\mathbb Z}na_n\e^{-\pi n^2x/N},
\]\[
\Lambda_{a,N}(s)=\frac{1}{2}\int_0^\infty\Theta_{a,N}(x)x^{(s+1)/2}\frac{dx}{x},
\]
the preceding residue-class proof, replacing $5$ by $N$, gives
\[
\Theta_{a,N}(x)=x^{-3/2}\Theta_{\Tcal_Na,N}(1/x),
\]\[
\Lambda_{a,N}(s)=\left(\frac{N}{\pi}\right)^{(s+1)/2}
\Gamma((s+1)/2)F_a(s)=\Lambda_{\Tcal_Na,N}(1-s).
\]
The series expression initially holds for $\Re s>1$. The integral is entire;
the same endpoint bounds and local-uniform differentiation argument apply.
This is a finite-Fourier theorem, distinct from an unconstructed
$N$-dependent surface.

\begin{thebibliography}{9}
\bibitem{Balanzario} E. P. Balanzario and J. Sanchez-Ortiz, \emph{Zeros of the Davenport--Heilbronn counterexample}, Mathematics of Computation \textbf{76} (2007), 2045--2049. doi:10.1090/S0025-5718-07-01999-0.
\bibitem{DLMF} NIST Digital Library of Mathematical Functions, \S1.8(iv), \url{https://dlmf.nist.gov/1.8#iv}.
\bibitem{Duke} W. Duke, \emph{Continued Fractions and Modular Functions}, Bulletin of the American Mathematical Society (N.S.) \textbf{42}(2) (2005), 137--162. doi:10.1090/S0273-0979-05-01047-5.
\bibitem{Hermite} C. Hermite, \emph{Sur la th\'eorie des \'equations modulaires et la r\'esolution de l'\'equation du cinqui\`eme degr\'e}, Paris: Mallet-Bachelier, 1859, pp.~47--48. \url{https://gdz.sub.uni-goettingen.de/download/pdf/PPN598770909/PPN598770909.pdf}.
\bibitem{HMR} A. Hayami, \emph{The Orthogonal-Circle Ruled Surface: A Boundary Atlas, Five Rank Singularities, and Conditional Fold-Signal Recurrence}, version 5, 2026, Definition~1.1 and Theorem~2.2. \href{https://doi.org/10.5281/zenodo.23073642}{doi:10.5281/zenodo.23073642}.
\bibitem{SpearmanWilliams} B. K. Spearman and K. S. Williams, \emph{The conductor of a cyclic quartic field using Gauss sums}, Czechoslovak Mathematical Journal \textbf{47} (1997), no.~3, 453--462. \url{https://dml.cz/handle/10338.dmlcz/127369}.
\end{thebibliography}
\end{document}
